\documentclass[11pt,a4paper]{article}
\usepackage[utf8]{inputenc}
\usepackage{amsmath}
\usepackage{graphicx}
\usepackage{booktabs}
\usepackage{hyperref}
\usepackage{geometry}
\geometry{margin=1in}

\title{
    \textbf{Cardiovascular Disease Risk Classification Pipeline}\\
    \large Capstone Project Report -- AI \& Data Science Basics (Semester S4)
}
\author{Assadiki Zaid\\ Academic Mentor: Dr. Rym Nassih}
\date{June 18, 2026}

\begin{document}

\maketitle
\thispagestyle{empty}
\newpage

\pagenumbering{roman}
\tableofcontents
\newpage

\pagenumbering{arabic}

\section{Executive Summary}
Early detection of cardiovascular complications can drastically alter patient outcomes and ease the critical burden on healthcare infrastructure. This project focuses on building a predictive pipeline to determine the presence or risk level of cardiovascular disease based on clinical measurements, lifestyle habits, and patient biometric data. By analyzing heterogeneous health indicators, the machine learning model serves as a clinical decision-support tool designed for medical professionals and healthcare providers. It enables them to identify high-risk individuals accurately, prioritize preventative therapeutic interventions, and optimize patient monitoring strategies before severe clinical events occur.

\section{Phase 1: Problem Framing \& AI Context Note}
The objective of this project is to develop a machine learning pipeline that categorizes patients into high-risk ($cardio=1$) or low-risk ($cardio=0$) based on routine clinical vitals, body measurements, and lifestyle surveys. 

\subsection{AI Context Note}
In early AI (Course 2, Chapters 1-3), the dominant paradigm was Symbolic AI and Rule-Based Expert Systems, such as the MYCIN system developed in the 1970s. These systems relied on hand-coded rules created by clinical experts:
\begin{equation}
\text{IF } \textit{systolic\_bp} > 140 \text{ AND } \textit{cholesterol} = 3 \text{ THEN } \textit{cardio\_risk} = \text{HIGH}
\end{equation}
While MYCIN and other expert systems were highly transparent, they suffered from the knowledge acquisition bottleneck, fragility, and an inability to learn from raw empirical data or manage high-dimensional probabilistic boundaries. The modern paradigm shift to empirical machine learning (connectionism) resolves this. Instead of manual rules, models like XGBoost learn associations directly from a dataset of 70,000 real-world observations, mapping continuous risk boundaries that are more robust to noisy data.

\section{Phase 2: Data Audit, Outlier Treatment, \& Encoding}
A comprehensive audit was performed on the raw Kaggle dataset (70,000 records). 
\subsection{Data Audit Table}
\begin{table}[htbp]
\centering
\caption{Data Quality Audit}
\begin{tabular}{llll}
\toprule
Variable & Type & Range & Audit Findings \\
\midrule
age & Int (days) & 10,798 -- 23,713 & Stored in days; requires conversion to years \\
gender & Cat (1/2) & 1 -- 2 & Re-encoded to binary (0/1) \\
height & Int (cm) & 55 -- 250 & Contains unphysical height values \\
weight & Float (kg) & 10 -- 200 & Contains unphysical weight values \\
ap\_hi & Int (mmHg) & -150 -- 16,000 & Severe outliers (logging errors) \\
ap\_lo & Int (mmHg) & -70 -- 11,000 & Severe outliers (logging errors) \\
cardio & Binary (0/1) & 0 -- 1 & Target variable; perfectly balanced \\
\bottomrule
\end{tabular}
\end{table}

\subsection{Outlier Comparison}
We compared the Interquartile Range (IQR) and Z-score methods:
\begin{itemize}
    \item \textbf{IQR Method:} Flagged 1,439 outliers.
    \item \textbf{Z-score Method ($|Z| > 3$):} Flagged 983 outliers.
    \item \textbf{Justification:} The Z-score method was heavily distorted by extreme values (e.g., blood pressure of 16,000 mmHg), which artificially inflated the standard deviation. A strict physical biomedical threshold was applied to clean the cohort, removing 6,058 records and leaving a clean sample of 63,942 patients.
\end{itemize}

\subsection{Categorical Encoding}
Ordinal features \textit{cholesterol} and \textit{gluc} were one-hot encoded to prevent the algorithms from assuming a strictly linear increment between normal (1), above normal (2), and well above normal (3).

\section{Phase 3: Feature Engineering \& Data Management}
We engineered three clinical indicators:
\begin{enumerate}
    \item \textbf{Body Mass Index (BMI):} $\text{BMI} = \text{weight (kg)} / (\text{height (m)})^2$
    \item \textbf{Pulse Pressure:} $\text{ap\_hi} - \text{ap\_lo}$ (arterial stiffness marker)
    \item \textbf{Age-Cholesterol Risk:} $\text{Age (Years)} \times \text{cholesterol}$
\end{enumerate}

The cleaned dataset was uploaded to a local SQLite database (\texttt{data/cardio\_cleaned.db}).

\section{Phase 4 \& 5: Modeling \& Performance Summary}
A train/test split (80/20) was executed with zero data leakage. We evaluated four models using pipelines:
\begin{table}[htbp]
\centering
\caption{Model Benchmarking Comparison}
\begin{tabular}{llllll}
\toprule
Model & Accuracy & Precision & Recall & $F_1$-score & ROC-AUC \\
\midrule
\textbf{Tuned XGBoost} & \textbf{73.57\%} & \textbf{75.39\%} & \textbf{70.21\%} & \textbf{72.71\%} & \textbf{80.12\%} \\
Logistic Regression & 72.84\% & 74.92\% & 68.61\% & 71.63\% & 79.28\% \\
Random Forest & 71.21\% & 72.10\% & 69.15\% & 70.59\% & 77.40\% \\
K-Nearest Neighbors & 69.80\% & 70.35\% & 68.10\% & 69.21\% & 74.50\% \\
\bottomrule
\end{tabular}
\end{table}

\subsection{Justification of $F_1$-Score}
In healthcare, raw accuracy is a misleading metric. We must minimize False Negatives (patients with active CVD diagnosed as healthy, which could be fatal) and False Positives (wasting clinical stress-test resources on healthy patients). The $F_1$-score balances Precision and Recall, protecting both the patient's health and the clinic's resources.

\section{Phase 6: Model Card \& Ethical Considerations}
\begin{itemize}
    \item \textbf{Intended Use:} Decision-support tool for clinic triage.
    \item \textbf{Operational Limits:} Not validated for pediatric patients ($<18$ years) or inputs outside human physiology.
    \item \textbf{Ethical Reflection:} Clinicians must oversee all automated recommendations (human-in-the-loop) and watch for gender-related diagnostic disparities.
\end{itemize}

\section{Phase 7: Disclosures \& Academic Integrity}
\subsection{AI-Assisted Content Appendix}
Generative AI was utilized to draft portions of this document and check programming structures. All outputs were verified by executing the scripts locally.

\subsection{Individual Contribution Statement}
As the sole author, I, Assadiki Zaid, am responsible for all components of this project.

\end{document}
